\section*{Bab \ref{ch:b1:diagenesis-basins} --- Diagenesis dan cekungan sedimen}

\begin{solution}{exo:b1:diagenesis-basins:1}
Batu pasir $0.49\,e^{-2000/3704} = 0.49\times0.583 = \num{0.286}$; serpih
$0.63\,e^{-2000/1960} = 0.63\times0.360 = \num{0.227}$. Porositas menjadi setengah
ketika $e^{-z/\lambda} = 1/2$, pada $z = \lambda\ln 2$: \qty{2567}{m} untuk batu pasir,
\qty{1359}{m} untuk serpih.
\end{solution}

\begin{solution}{exo:b1:diagenesis-basins:2}
$\Delta A = \Delta E + S = 30 + 50 = \qty{80}{m}$. Darinya, \qty{60}{m} terisi dengan
sedimen, sehingga $\Delta W = 80 - 60 = \qty{20}{m}$: paparan menjadi lebih dalam
sebesar \qty{20}{m}.
\end{solution}

\begin{solution}{exo:b1:diagenesis-basins:3}
$h = Y(\rho_m - \rho_w)/(\rho_m - \rho_s)$: dengan \qty{2300}{kg.m^{-3}},
$1.0\times2270/1000 = \qty{2.27}{km}$; dengan \qty{2500}{kg.m^{-3}},
$1.0\times2270/800 = \qty{2.84}{km}$. Setiap meter sedimen menempati tempat satu meter
\omterm{def:b1:layered-earth:layers}{mantel} dalam kolom; makin padat sedimennya, makin kecil massa yang kurang dibanding
\omterm{def:b1:layered-earth:layers}{mantel}, sehingga makin banyak yang dibutuhkan untuk mengisi defisit yang sama dengan
lubang berisi air.
\end{solution}

\begin{solution}{exo:b1:diagenesis-basins:4}
(a)~Kalsit $100.09/2.71 = \qty{36.93}{cm^3.mol^{-1}}$; dolomit $184.40/2.84
= \qty{64.93}{cm^3.mol^{-1}}$. (b)~Dua mol kalsit menempati \qty{73.86}{cm^3}, satu mol
dolomit \qty{64.93}{cm^3}: padatan menyusut sebesar $1 - 64.93/73.86 = \qty{12.1}{\%}$.
(c)~Dalam satu meter kubik, \qty{0.90}{m^3} kalsit menjadi $0.90\times(1 - 0.121) = \qty{0.791}{m^3}$ dolomit; porositasnya menjadi $1 - 0.791 = \num{0.21}$.
\end{solution}

\begin{solution}{exo:b1:diagenesis-basins:5}
(a)~$s = 1000 - 0.63\times1960\times(1 - e^{-1000/1960}) = \qty{506.5}{m}$; porositas
rerata adalah $1 - 506.5/1000 = \num{0.493}$. (b)~Dengan puncak di \qty{2000}{m},
menyelesaikan $s(2000, b) = 506.5$ memberi $b = \qty{2628.7}{m}$: lapisan itu setebal
\qty{628.7}{m}, dengan porositas rerata $1 - 506.5/628.7 = \num{0.194}$. (c)~Butirannya
sama, sehingga air yang hilang sama dengan ketebalan yang hilang: $1000 - 628.7 = \qty{371}{m^3}$ air per meter persegi.
\end{solution}

\begin{solution}{exo:b1:diagenesis-basins:6}
(a)~$s = 300 - 0.49\times3704\times(e^{-2500/3704} - e^{-2800/3704}) =
\qty{228.10}{m}$. (b)~Dari $b_0 = 300$: $s(0, 300) = 158.80$, $1 - \varphi(300) =
0.5481$, $b_1 = 300 + (228.10 - 158.80)/0.5481 = \qty{426.44}{m}$; lalu $s(0, 426.44) = 229.07$, $1 - \varphi = 0.5633$, $b_2 = 426.44 - 0.96/0.5633 =
\qty{424.73}{m}$; langkah ketiga mengubahnya kurang dari satu milimeter. Batu pasir itu
setebal \qty{424.7}{m} di permukaan. (c)~$424.7/300 = 1.42$: penguburan telah
menipiskannya dengan faktor \num{1.42}.
\end{solution}

\begin{solution}{exo:b1:diagenesis-basins:7}
(a)~Es memegang $52\times10^6\times916.7/1000 = \qty{4.77e7}{km^3}$ air; tersebar di
atas \qty{361.84e6}{km^2} ia adalah $\qty{0.1317}{km} = \qty{132}{m}$ muka laut.
(b)~Garis pantai memberi \qty{134}{m}, \qty{2}{m} lebih. (c)~$4.767\times
10^7/0.134 = \qty{355.7e6}{km^2}$, \qty{1.7}{\%} lebih kecil dari hari ini. Ketika muka
laut turun, paparan benua muncul, sehingga samudra glasial lebih kecil, dan volume air
yang sama yang diambil darinya menurunkannya lebih banyak.
\end{solution}

\begin{solution}{exo:b1:diagenesis-basins:8}
(a)~$(-9.1 - (-98.1))/8 = \qty{11.1}{m}$ per seribu tahun, mendekati \qty{12}{m} yang
ditemukan dari garis pantai. (b)~\qty{40}{mm} setahun adalah \qty{40}{m} per seribu
tahun, $40/12 = 3.3$ kali rerata; dalam lima ratus tahun, $0.040\times500 = \qty{20}{m}$.
\end{solution}

\begin{solution}{exo:b1:diagenesis-basins:9}
(a)~$S = 50\times2 = \qty{100}{m}$ dan $\Delta E = \qty{-40}{m}$: $\Delta A = \qty{60}{m}$, seluruhnya terisi sedimen karena kedalaman tidak berubah. (b)~Dalam
\num{20000} tahun batuan dasar turun hanya $50\times0.02 = \qty{1}{m}$: \omterm{def:b1:diagenesis-basins:eustasy}{muka laut
relatif} turun sebesar $120 - 1 = \qty{119}{m}$. Dasar laut, yang berada \qty{20}{m}
di bawah muka laut, berakhir $119 - 20 = \qty{99}{m}$ di atasnya: seluruh paparan muncul.
\end{solution}

\begin{solution}{exo:b1:diagenesis-basins:10}
(a)~Untuk lapisan tipis porositasnya hampir konstan di seluruhnya: ketebalan padatnya
adalah $s \approx h(1 - \varphi(z))$ pada kedalaman dan $s \approx h_0(1 -
\varphi_0)$ di permukaan. Menyamakan keduanya memberi $h_0 \approx h(1 -
\varphi(z))/(1 - \varphi_0)$. (b)~$\varphi(3000) = 0.63\,e^{-3000/1960} =
0.136$, sehingga $h_0 \approx 10\times0.864/0.37 = \qty{23.3}{m}$, dibandingkan
\qty{23.1}{m} secara tepat. (c)~Dipulihkan ke permukaan, lapisan itu setebal
\qty{23}{m}, tidak lagi tipis dibandingkan dengan peluruhan porositas di dekat
permukaan: porositasnya turun dari $\varphi_0$ di puncaknya menjadi lebih kecil di
dasarnya, sehingga butirannya muat dalam ketebalan yang sedikit lebih kecil daripada
yang dihitung dengan $\varphi_0$ di seluruh lapisan.
\end{solution}

\begin{solution}{exo:b1:diagenesis-basins:11}
(a)~$s = 3000 - 0.63\times1960\times(1 - e^{-3000/1960}) = \qty{2032}{m}$; porositas
rerata adalah $1 - 2032/3000 = \num{0.323}$ dan massa jenis rerata
$0.323\times1030 + 0.677\times2700 = \qty{2161}{kg.m^{-3}}$. (b)~Faktornya adalah
$(3300 - 1030)/(3300 - 2161) = \num{1.99}$; isian itu mewakili \omterm{def:b1:diagenesis-basins:subsidence}{penurunan} beban-air
sebesar $3000/1.99 = \qty{1505}{m}$. (c)~Dengan massa jenis butiran faktornya akan
$2270/600 = \num{3.78}$ dan \omterm{def:b1:diagenesis-basins:subsidence}{penurunannya} hanya \qty{793}{m}. Itu keliru karena sepertiga
isian adalah air: sedimen jauh lebih ringan daripada butirannya, dan hanya massa jenis sedimen
jenuh yang dihitung dalam keseimbangan massa.
\end{solution}

\begin{solution}{exo:b1:diagenesis-basins:12}
(a)~Kubus berisi \qty{0.90}{m^3} $= \qty{9.0e5}{cm^3}$ kalsit, yakni $9.0\times10^5/36.93 = \num{24370}$ mol. (b)~Satu magnesium per dua kalsit: \num{12180} mol. (c)~Satu
kilogram air laut mengandung $35.165\times0.0365 =
\qty{1.284}{g}$ magnesium,
$1.284/24.305 = \qty{0.0528}{mol}$; reaksi itu membutuhkan $12180/0.0528 = \qty{2.31e5}{kg}$ air laut, $\qty{224}{m^3}$ --- lebih dari dua ratus kali volume batuan,
bahkan jika setiap ion magnesium diambil. (d)~Dolomit dapat terbentuk dalam skala besar
hanya di tempat volume air yang sangat besar mengalir melalui batuan dalam waktu lama,
atau di tempat air itu adalah air asin yang lebih kaya magnesium daripada air laut:
misalnya di platform yang dibilas oleh air asin yang tenggelam, dan bukan di sedimen
yang mempertahankan air porinya yang sama.
\end{solution}

\begin{solution}{pb:b1:diagenesis-basins:1}
\textbf{1.} Batu pasir: $\varphi(0) = \num{0.490}$, $\varphi(800) =
0.49\,e^{-800/3704} = \num{0.395}$. Serpih: $\varphi(800) = \num{0.419}$,
$\varphi(2600) = \num{0.167}$. Kapur: $\varphi(2600) = \num{0.110}$,
$\varphi(3200) = \num{0.072}$.

\textbf{2.} Dengan $s$ dari \cref{prop:b1:diagenesis-basins:solid} dan $\bar\varphi
= 1 - s/h$ (soal 7): batu pasir \num{0.441}, serpih \num{0.274}, kapur
\num{0.090}.

\textbf{3.} $h\bar\varphi$: batu pasir \qty{353}{m^3}, serpih \qty{493}{m^3}, kapur
\qty{54}{m^3} air per meter persegi.

\textbf{4.} $\bar\varphi\rho_w + (1 - \bar\varphi)\rho_g$: batu pasir
$0.441\times1030 + 0.559\times2650 = \qty{1936}{kg.m^{-3}}$, serpih
\qty{2242}{kg.m^{-3}}, kapur \qty{2559}{kg.m^{-3}}. Kolom itu berbobot
$1936\times800 + 2242\times1800 + 2559\times600 = \qty{7.12e6}{kg}$ per meter persegi:
massa jenis rerata \qty{2225}{kg.m^{-3}}.

\textbf{5.} Litostatik $2225\times9.81\times3200 = \qty{69.9}{MPa}$; hidrostatik
$1030\times9.81\times3200 = \qty{32.3}{MPa}$; \omterm{def:b1:stress-strain-rheology:lithostatic}{tegangan efektif} $69.9 - 32.3 = \qty{37.5}{MPa}$.

\textbf{6.} $0.49\,e^{-z/3704} = 0.63\,e^{-z/1960}$ memberi $z = \ln(0.63/0.49)
/(1/1960 - 1/3704) = \qty{1046}{m}$. Di bawahnya batu pasir lebih berpori: serpih
mulai lebih basah tetapi kehilangan airnya dalam jarak yang lebih pendek.

\textbf{7.} Batu pasir $800 - 0.49\times3704\times(1 - e^{-800/3704}) =
\qty{447.4}{m}$; serpih \qty{1306.7}{m}; kapur \qty{546.0}{m}.

\textbf{8.} \qty{2300.2}{m} butiran; air membentuk $1 - 2300.2/3200 =
\qty{28}{\%}$ dari kolom.

\textbf{9.} Batu pasir $447.4/15 = \qty{29.8}{m}$, serpih $1306.7/40 =
\qty{32.7}{m}$, kapur $546.0/20 = \qty{27.3}{m}$ padatan per juta tahun.

\textbf{10.} Serpih memiliki pasokan tercepat, kapur yang paling lambat --- tetapi
ketiganya berbeda hanya sekitar seperlima, jauh lebih sedikit daripada yang disarankan
oleh ketebalan mereka hari ini.

\textbf{11.} \qty{800}{m}: puncaknya berada di permukaan hari ini, tempat pengendapannya
baru saja berakhir, sehingga ketebalannya saat diendapkan sama dengan ketebalan hari ini.

\textbf{12.} Menyelesaikan $s(0, b) = 1306.7$ untuk serpih dengan
\cref{met:b1:diagenesis-basins:decompaction} memberi $b = \qty{2124}{m}$: serpih setebal
\qty{2124}{m} pada \qty{15}{\Ma}.

\textbf{13.} $s(0, b) = 546.0$ untuk kapur memberi \qty{1071}{m} pada \qty{55}{\Ma}.

\textbf{14.} Dengan puncaknya di \qty{2124}{m}, $s(2124, b) = 546.0$ memberi $b =
\qty{2748}{m}$: kapur setebal \qty{624}{m}.

\textbf{15.} \qty{1071}{m} pada \qty{55}{\Ma}; $2124 + 624 = \qty{2748}{m}$ pada
\qty{15}{\Ma}.

\textbf{16.} Dari ketebalan hari ini: batu pasir $800/15 = \qty{53.3}{m}$, serpih
$1800/40 = \qty{45.0}{m}$, kapur $600/20 = \qty{30.0}{m}$ per juta tahun. Dari
ketebalan saat diendapkan: \qty{53.3}{m}, $2124/40 =
\qty{53.1}{m}$ dan $1071/20 = \qty{53.6}{m}$. Himpunan kedua adalah yang bermakna: laju tempat dasar laut dibangun,
dan ia menunjukkan laju yang tetap yang disembunyikan oleh ketebalan hari ini.

\textbf{17.} $1 - 1800/2124 = \qty{15}{\%}$.

\textbf{18.} Kedalaman air sama pada awal dan akhir setiap interval, sehingga \omterm{def:b1:diagenesis-basins:accommodation}{akomodasi}
yang tercipta adalah pertumbuhan kolom di atas datum: \qty{1071}{m} dari 75 hingga
\qty{55}{\Ma}, $2748 - 1071 = \qty{1677}{m}$ dari 55 hingga \qty{15}{\Ma}, dan
$3200 - 2748 = \qty{452}{m}$ sejak \qty{15}{\Ma}.

\textbf{19.} $S = \Delta A - \Delta E$: $1677 - (-60) = \qty{1737}{m}$ dari 55 hingga
\qty{15}{\Ma}; $452 - (-30) = \qty{482}{m}$ sejak itu.

\textbf{20.} $1737/40 = \qty{43}{m}$ dan $482/15 = \qty{32}{m}$ per juta tahun.

\textbf{21.} $60/1737 = \qty{3.5}{\%}$ dan $30/482 = \qty{6.2}{\%}$. Muka laut turun pada
kedua interval, sehingga \omterm{def:b1:diagenesis-basins:eustasy}{eustasi} menghambat pengisian: dasar harus turun sebanyak itu
lagi untuk menciptakan ruang yang sama. \omterm{def:b1:diagenesis-basins:subsidence}{Penurunan} menghasilkan hampir seluruh \omterm{def:b1:diagenesis-basins:accommodation}{akomodasi}.

\textbf{22.} Serpih dan kapur setebal \qty{2748}{m} pada \qty{15}{\Ma} dan kini
\qty{2400}{m}: \omterm{def:b1:diagenesis-basins:compaction}{pemadatan} mereka menciptakan \qty{348}{m} ruang, \qty{77}{\%} dari
\qty{452}{m} yang diisi oleh batu pasir.

\textbf{23.} Setiap satuan telah menipis oleh penguburan yang mengikutinya, dengan jumlah
yang bergantung pada litologinya dan seberapa dalam ia terkubur; ruang yang tercipta di
setiap interval harus diukur dengan ketebalan yang dimiliki satuan-satuan itu saat itu,
yang hanya dapat diberikan oleh \omterm{def:b1:diagenesis-basins:compaction}{dekompaksi}. Dan bahkan kemudian \omterm{def:b1:diagenesis-basins:subsidence}{penurunan} yang ditemukan
dengan cara itu mencakup \omterm{def:b1:diagenesis-basins:subsidence}{penurunan} yang disebabkan oleh berat sedimen, yang harus
dihilangkan untuk mencapai \omterm{def:b1:diagenesis-basins:subsidence}{penurunan tektonik}.

\textbf{24.} \qty{2124}{m}, sekitar \qty{2.1}{km} --- \qty{324}{m} lebih banyak daripada
\qty{1800}{m} yang ditemukan sumur hari ini.
\end{solution}
